\documentclass[10pt,twocolumn,letterpaper]{article}

% --- PREAMBLE: Packages & Configurations ---
\usepackage{times}
\usepackage{graphicx}
\usepackage{amsmath}
\usepackage{amssymb}
\usepackage{booktabs}
\usepackage{cite}
\usepackage[pagebackref,breaklinks,colorlinks]{hyperref}

% --- METADATA: Title & Authors ---
\title{Title of Research Paper}
\author{Author Name\\Institution}
\date{}

\begin{document}

\maketitle

\begin{abstract}
We study self-supervised clip ranking for autonomous driving video, with the goal of assigning a review-value score to short clips before human inspection. The paper evaluates a JEPA-style prediction objective as a proxy for novelty and review priority, and reports both ranking quality and inference efficiency. We use a three-way split protocol: \texttt{train} for fitting, \texttt{val} for checkpoint selection and development-time analysis, and \texttt{test} for final reported results.
\end{abstract}

\section{Introduction}
This paper addresses the problem of ranking driving clips by likely review value. In large-scale autonomous vehicle datasets, the majority of clips are routine, while only a small subset contains rare, informative, or failure-relevant events. A useful ranking model should surface those clips early while remaining computationally practical.

Our working hypothesis is that a JEPA-style representation predictor can produce a meaningful unsupervised novelty signal for clip ranking. The core question is whether this proxy improves prioritization quality enough to justify its inference cost.

\section{Related Work}
This section should position the work relative to:
\begin{itemize}
    \item self-supervised video representation learning,
    \item masked video modeling and JEPA-style prediction objectives,
    \item active curation, novelty detection, and clip triage for driving data,
    \item ranking metrics for retrieval and prioritization.
\end{itemize}

\section{Problem Setup}
Let each clip be a short temporal sequence of frames with optional human review labels available only for benchmarking. The model produces a scalar \texttt{review\_value\_score} for each clip, and evaluation measures whether high-scoring clips align with high-value human judgments.

The paper should define:
\begin{itemize}
    \item the clip construction process,
    \item the label taxonomy used for benchmarking,
    \item the ranking task and evaluation targets,
    \item the distinction between unsupervised scoring and supervised benchmarking.
\end{itemize}

\section{Method}
\subsection{Architectural Overview}
Describe the encoder, predictor, masking strategy, and tubelet-based processing pipeline.

\subsection{Training Objective}
Describe the embedding prediction objective and any normalization or loss choices used during training.

\subsection{Scoring Rule}
Describe how clip-level scores are aggregated from tubelet-level outputs. In the current implementation, the score is derived from prediction novelty via cosine disagreement between predicted and target embeddings.

\section{Experimental Setup}
\subsection{Dataset and Splits}
Use the split protocol explicitly:
\begin{itemize}
    \item \texttt{train}: fit model parameters,
    \item \texttt{val}: select checkpoints, tune settings, and analyze ranking behavior during development,
    \item \texttt{test}: report final paper metrics once the method is fixed.
\end{itemize}

This split usage should be stated clearly to avoid mixing development-time decisions with final reported results.

\subsection{Implementation Details}
Document the model variant, initialization strategy, encoder freeze/finetune setting, batch sizes, hardware profile, and runtime configuration used in experiments.

\section{Results on Validation}
This section should summarize development-time findings:
\begin{itemize}
    \item validation ranking metrics used to compare variants,
    \item checkpoint selection criteria,
    \item ablations on initialization, model capacity, or scoring aggregation,
    \item qualitative inspection of highly ranked clips.
\end{itemize}

\section{Final Test Evaluation}
This section should contain the headline paper numbers. Once model choices are fixed on \texttt{val}, run the final scoring and labeled evaluation on \texttt{test} and report the primary ranking metrics here.

\section{Efficiency and Inference Cost}
Report throughput, latency, peak memory, and estimated energy usage for scoring. This section should make the tradeoff between ranking quality and deployment cost explicit.

\section{Discussion and Limitations}
Discuss what the novelty proxy captures well, where it fails, how sensitive it is to split construction and label quality, and whether ranking quality generalizes from \texttt{val} to \texttt{test}.

\section{Conclusion}
Summarize the main empirical findings, the practical value of the ranking pipeline, and the most important next steps for improving both quality and efficiency.

\section*{Acknowledgments}
Optional.

{\small
\bibliographystyle{ieee_fullname}
\bibliography{references}
}

\end{document}
